\documentclass[12pt,oneside]{book}

% ============================================================
% METADATOS
% ============================================================
\newcommand{\universidad}{Universidad de Buenos Aires (UBA)}
\newcommand{\materia}{Derecho Procesal Civil}
\newcommand{\titulo}{Apuntes de Derecho Procesal Civil}
\newcommand{\subtitulo}{Recursos procesales y Rebeld\'ia}
\newcommand{\autor}{Isaac Edgar Camacho Ocampo}
\newcommand{\anio}{2026}

% ============================================================
% PAQUETES
% ============================================================
\usepackage[utf8]{inputenc}
\usepackage[T1]{fontenc}
\usepackage[english]{babel}
\addto\captionsenglish{\renewcommand{\contentsname}{Índice}}
\usepackage[
    top=2.5cm,
    bottom=2.5cm,
    left=3cm,
    right=2.5cm,
    headheight=15pt
]{geometry}
\usepackage{amsmath}
\usepackage{amssymb}
\usepackage{mathtools}
\usepackage{graphicx}
\usepackage{float}
\usepackage{caption}
\usepackage{subcaption}
\usepackage{booktabs}
\usepackage{array}
\usepackage{longtable}
\usepackage{tabularx}
\usepackage{enumitem}
\usepackage{xcolor}
\usepackage[most]{tcolorbox}
\usepackage{fancyhdr}
\usepackage{ragged2e}
\usepackage[
    colorlinks=true,
    linkcolor=blue!50!black,
    citecolor=green!40!black,
    urlcolor=blue!60!black,
    pdftitle={\titulo},
    pdfauthor={\autor},
    pdfsubject={Apuntes universitarios},
    bookmarks=true
]{hyperref}

% ============================================================
% CONFIGURACIÓN
% ============================================================
\graphicspath{
    {./img/}
    {./graficos/}
    {./figuras/}
}
\setcounter{tocdepth}{2}
\setlength{\parindent}{0pt}
\setlength{\parskip}{0.5em}
\setlist[itemize]{
    topsep=0.3em,
    itemsep=0.2em
}
\setlist[enumerate]{
    topsep=0.3em,
    itemsep=0.2em
}
\clubpenalty=10000
\widowpenalty=10000

% ============================================================
% COMANDOS
% ============================================================
\newcommand{\abs}[1]{\left|#1\right|}
\newcommand{\norm}[1]{\left\lVert#1\right\rVert}
\newcommand{\vect}[1]{\mathbf{#1}}
\newcommand{\deriv}[2]{\frac{d#1}{d#2}}
\newcommand{\pderiv}[2]{\frac{\partial#1}{\partial#2}}

% ============================================================
% ENTORNOS / CAJAS
% ============================================================
\newtcolorbox{definition}[1]{
    enhanced,
    breakable,
    title={Definici\'on: #1},
    fonttitle=\bfseries,
    colback=blue!3,
    colframe=blue!50!black,
    boxrule=0.8pt,
    arc=2mm
}

\newtcolorbox{important}{
    enhanced,
    breakable,
    title={Importante},
    fonttitle=\bfseries,
    colback=yellow!8,
    colframe=orange!70!black,
    boxrule=0.8pt,
    arc=2mm
}

\newtcolorbox{example}[1]{
    enhanced,
    breakable,
    title={Ejemplo: #1},
    fonttitle=\bfseries,
    colback=green!4,
    colframe=green!40!black,
    boxrule=0.8pt,
    arc=2mm
}

\newtcolorbox{formula}[1]{
    enhanced,
    breakable,
    title={F\'ormula / Regla: #1},
    fonttitle=\bfseries,
    colback=purple!4,
    colframe=purple!50!black,
    boxrule=0.8pt,
    arc=2mm
}

\newtcolorbox{exercise}[1]{
    enhanced,
    breakable,
    title={Ejercicio: #1},
    fonttitle=\bfseries,
    colback=gray!5,
    colframe=gray!60!black,
    boxrule=0.8pt,
    arc=2mm
}

\newtcolorbox{note}{
    enhanced,
    breakable,
    title={Nota},
    fonttitle=\bfseries,
    colback=white,
    colframe=gray!50,
    boxrule=0.6pt,
    arc=2mm
}

% ============================================================
% CONFIGURACIÓN FINAL
% ============================================================
\pagestyle{fancy}
\fancyhf{}
\fancyhead[L]{\nouppercase{\leftmark}}
\fancyhead[R]{\materia}
\fancyfoot[C]{\thepage}
\renewcommand{\headrulewidth}{0.4pt}
\renewcommand{\footrulewidth}{0pt}

\fancypagestyle{plain}{
    \fancyhf{}
    \fancyfoot[C]{\thepage}
    \renewcommand{\headrulewidth}{0pt}
}

\begin{document}

% ============================================================
% PORTADA
% ============================================================
\begin{titlepage}

\thispagestyle{empty}

\begin{center}

\vspace*{1cm}

{\large \textsc{\universidad}}

\vspace{3cm}

{\Huge \textbf{\materia}}

\vspace{1cm}

{\LARGE \subtitulo}

\vspace{0.8cm}

\textit{Comprender los conceptos es el primer paso para convertir el estudio en conocimiento.}

\vspace{1.2cm}

\rule{0.70\textwidth}{0.5pt}

\vspace{0.5cm}

{\large Apuntes de estudio}

\vspace{0.5cm}

\rule{0.70\textwidth}{0.5pt}

\vfill

{\large \autor}

\vspace{0.3cm}

{\large \anio}

\end{center}

\vfill

\begin{flushleft}
\textit{Este es un modesto aporte a los estudiantes de ``Derecho Procesal Civil'' no pretende sustituir el material oficial de la materia.}
\end{flushleft}

\end{titlepage}

% ============================================================
% ÍNDICE
% ============================================================
\tableofcontents

% ============================================================
% CONTENIDO
% ============================================================

% ------------------------------------------------------------
\chapter[Los recursos procesales]{Los recursos procesales}
% ------------------------------------------------------------

\section{Concepto de recurso}

\begin{definition}{Recurso}
El \textbf{recurso} es un \textbf{medio procesal} mediante el cual la parte que se siente agraviada por una resoluci\'on judicial puede solicitar la \textbf{revocaci\'on total o parcial} de esa resoluci\'on, ante el mismo juez que la dict\'o o ante otro de jerarqu\'ia superior.
\end{definition}

\section{Clasificaci\'on y reglas b\'asicas}

\subsection{Recursos ordinarios y extraordinarios}

La ley clasifica los recursos en \textbf{ordinarios} y \textbf{extraordinarios}. Los ordinarios son aquellos que se conceden en cualquier etapa del proceso. Los extraordinarios se conceden de manera excepcional y su procedimiento est\'a espec\'ificamente determinado por la ley.

Dentro de los recursos ordinarios se encuentran la aclaratoria, la reposici\'on o revocatoria, la apelaci\'on y la queja por apelaci\'on denegada. Dentro de los extraordinarios se encuentran el recurso extraordinario ante la Corte Suprema de Justicia de la Naci\'on y el del art\'iculo 14 de la Ley 48.

% TODO: contenido incompleto o ambiguo en las notas originales (las notas de edici\'on indican que ``art\'iculo 14 de la Ley 48'' corrige una inversi\'on de n\'umeros en el audio)

\begin{table}[H]
\centering
\renewcommand{\arraystretch}{1.3}
\begin{tabularx}{\textwidth}{
    >{\raggedright\arraybackslash}p{0.2\textwidth}
    >{\raggedright\arraybackslash}X
    >{\raggedright\arraybackslash}X
}
\toprule
\textbf{Clasificaci\'on} & \textbf{Cu\'ando se concede} & \textbf{Recursos comprendidos} \\
\midrule
Ordinarios & En cualquier etapa del proceso & Aclaratoria; reposici\'on o revocatoria; apelaci\'on; queja por apelaci\'on denegada \\
Extraordinarios & Excepcionalmente, con el procedimiento espec\'ificamente determinado por la ley & Recurso extraordinario ante la Corte Suprema de Justicia de la Naci\'on (art. 14, Ley 48) \\
\bottomrule
\end{tabularx}
\end{table}

\subsection{Los tres puntos b\'asicos para interponer un recurso}

Para interponer cualquier recurso deben tenerse claros tres puntos:

\begin{enumerate}
    \item Saber \textbf{qu\'e recurso se interpone}.
    \item Conocer el \textbf{plazo}.
    \item Saber \textbf{c\'omo se fundamenta}.
\end{enumerate}

\begin{important}
Los tres puntos son: qu\'e se interpone, en qu\'e plazo y con qu\'e fundamentaci\'on. Teniendo esto claro, la interposici\'on de un recurso resulta sencilla.
\end{important}

\section{Recurso de aclaratoria}

Por medio de este recurso se pide al juez que dict\'o una resoluci\'on que \textbf{aclare alg\'un concepto}, por alguna omisi\'on o alg\'un error. Se interpone \textbf{ante el mismo juez} que dict\'o la resoluci\'on, en el plazo de \textbf{3 d\'ias}, y se fundamenta en el mismo escrito de interposici\'on.

Este recurso \textbf{no requiere sustanciaci\'on}: no se corre traslado a la parte contraria. Simplemente se solicita que se corrija un error, que se aclare una fecha o que se rectifique un nombre equivocado. Es el m\'as sencillo de todos los recursos.

\section{Recurso de revocatoria (reposici\'on)}

Por el recurso de revocatoria tambi\'en se pide al \textbf{mismo juez} que dict\'o la resoluci\'on que la revoque. Procede \textbf{contra las providencias simples \'unicamente}.

Debe interponerse dentro del plazo de \textbf{3 d\'ias} de notificada la resoluci\'on, se fundamenta en el mismo escrito y es fundamental que vaya acompa\~nado por la \textbf{apelaci\'on subsidiaria}.

\begin{important}
La revocatoria \textbf{no suspende el plazo de la apelaci\'on}. Por eso debe interponerse siempre con apelaci\'on subsidiaria: si se interpone solo la revocatoria y luego se la deniega, el plazo para apelar ya se ha perdido.
\end{important}

\section{Recurso de apelaci\'on}

\subsection{Concepto, plazo y procedencia}

Por medio del recurso de apelaci\'on se pide al \textbf{\'organo superior} que revoque la decisi\'on del inferior. Se interpone \textbf{ante el mismo juez} que dict\'o la resoluci\'on, dentro del plazo de \textbf{5 d\'ias}.

Procede contra:
\begin{itemize}
    \item las sentencias definitivas en juicios ordinarios;
    \item las sentencias interlocutorias;
    \item las providencias simples que causen gravamen irreparable (no cualquier providencia).
\end{itemize}

\subsection{Efectos del recurso}

El recurso de apelaci\'on tiene dos efectos: \textbf{suspensivo} y \textbf{devolutivo}. Se concede siempre con efecto suspensivo, salvo que la ley disponga lo contrario; ello significa que se suspende la sentencia. El efecto devolutivo, por el contrario, significa que la sentencia debe cumplirse aunque haya sido apelada.

Se concede con efecto devolutivo en las medidas cautelares, en las sentencias de alimentos y en las sentencias de los procesos sumar\'isimos. En esos casos siempre procede con efecto devolutivo.

Adem\'as, existe un efecto llamado \textbf{diferido}, que se vincula con el momento en que se fundamenta el recurso.

% TODO: contenido incompleto o ambiguo en las notas originales (el efecto diferido se menciona sin mayor desarrollo)

\begin{table}[H]
\centering
\renewcommand{\arraystretch}{1.3}
\begin{tabularx}{\textwidth}{
    >{\raggedright\arraybackslash}p{0.17\textwidth}
    >{\raggedright\arraybackslash}X
    >{\raggedright\arraybackslash}X
}
\toprule
\textbf{Efecto} & \textbf{Significado} & \textbf{Cu\'ando aplica} \\
\midrule
Suspensivo & Se suspende la sentencia & Regla general: siempre, salvo que la ley disponga lo contrario \\
Devolutivo & La sentencia se cumple aunque haya sido apelada & Medidas cautelares; sentencias de alimentos; sentencias de procesos sumar\'isimos \\
Diferido & Determina cu\'ando se fundamenta el recurso & Se vincula con el momento de la fundamentaci\'on \\
\bottomrule
\end{tabularx}
\end{table}

\subsection{Formas de concesi\'on: libremente o en relaci\'on}

El recurso se concede de dos formas: \textbf{libremente} o \textbf{en relaci\'on}. La apelaci\'on se concede libremente en la sentencia definitiva de los juicios ordinarios, y se concede en relaci\'on cuando se trata de sentencias interlocutorias o de providencias simples que causen gravamen irreparable.

\subsection{Fundamentaci\'on cuando se concede libremente: expresi\'on de agravios}

Cuando el recurso se concedi\'o libremente (contra sentencias definitivas en juicios ordinarios), la fundamentaci\'on se denomina \textbf{expresi\'on de agravios}. Presentada la apelaci\'on y concedido el recurso libremente, el expediente se eleva a la C\'amara de Apelaciones. Una vez que llega a la C\'amara, se notifica cu\'al es la sala que interviene y, a partir de ese momento, corre un plazo de \textbf{10 d\'ias} para acompa\~nar la expresi\'on de agravios.

Si en ese plazo no se presenta la expresi\'on de agravios, o no se realiza una cr\'itica razonada y fundada de la sentencia, el recurso se tiene por \textbf{desierto}.

\subsection{Fundamentaci\'on cuando se concede en relaci\'on: memorial}

Cuando el recurso se concede en relaci\'on, la fundamentaci\'on se denomina \textbf{memorial} y se presenta \textbf{ante el mismo juez} que dict\'o la resoluci\'on, dentro del plazo de \textbf{5 d\'ias}. Una vez notificada la providencia que tiene por presentado el memorial, se corre traslado a la parte contraria por el plazo de \textbf{5 d\'ias}. Cuando la parte contraria contest\'o el memorial, o venci\'o el plazo sin que lo contestara, se pide la elevaci\'on a la C\'amara.

\begin{table}[H]
\centering
\renewcommand{\arraystretch}{1.3}
\begin{tabularx}{\textwidth}{
    >{\raggedright\arraybackslash}p{0.2\textwidth}
    >{\raggedright\arraybackslash}X
    >{\raggedright\arraybackslash}X
}
\toprule
 & \textbf{Concedido libremente} & \textbf{Concedido en relaci\'on} \\
\midrule
Procede contra & Sentencia definitiva de juicios ordinarios & Sentencias interlocutorias y providencias simples que causen gravamen irreparable \\
Fundamentaci\'on & Expresi\'on de agravios & Memorial \\
Ante qui\'en se presenta & C\'amara de Apelaciones (la sala que interviene) & El mismo juez que dict\'o la resoluci\'on \\
Plazo & 10 d\'ias desde la notificaci\'on de la sala & 5 d\'ias \\
Traslado a la contraria & No se menciona & 5 d\'ias; luego se pide la elevaci\'on a la C\'amara \\
Si no se fundamenta & Se tiene por desierto & No se menciona \\
\bottomrule
\end{tabularx}
\end{table}

\subsection{Incidentes del art\'iculo 250 y copias}

Cuando el recurso de apelaci\'on se concede en relaci\'on se forman los llamados \textbf{incidentes del art\'iculo 250}. Para la formaci\'on del incidente deben acompa\~narse \textbf{copias}. Al conceder el recurso, el juzgado indica cu\'ales son las copias que deben acompa\~narse.

La raz\'on es que, si se concedi\'o el recurso en relaci\'on, la resoluci\'on apelada debe seguir cumpli\'endose:

\begin{itemize}
    \item \textbf{Sentencia definitiva} (por ejemplo, una sentencia de alimentos): el \textbf{expediente principal} va a la C\'amara y quedan en el juzgado las copias indicadas por \'este, para que la sentencia pueda seguir ejecut\'andose.
    \item \textbf{Sentencia interlocutoria} (por ejemplo, alimentos provisorios): el \textbf{expediente principal queda en el juzgado} y se remiten a la C\'amara el escrito de interposici\'on del recurso, los memoriales y las copias que el juzgado indique, para la formaci\'on del incidente en C\'amara.
\end{itemize}

\begin{table}[H]
\centering
\renewcommand{\arraystretch}{1.3}
\begin{tabularx}{\textwidth}{
    >{\raggedright\arraybackslash}p{0.24\textwidth}
    >{\raggedright\arraybackslash}X
    >{\raggedright\arraybackslash}X
}
\toprule
\textbf{Resoluci\'on apelada (en relaci\'on)} & \textbf{Qu\'e va a la C\'amara} & \textbf{Qu\'e queda en el juzgado} \\
\midrule
Sentencia definitiva (ej.: sentencia de alimentos) & El expediente principal & Las copias indicadas por el juzgado, para que la sentencia siga cumpli\'endose \\
Sentencia interlocutoria (ej.: alimentos provisorios) & El escrito de interposici\'on del recurso, los memoriales y las copias que indique el juzgado & El expediente principal \\
\bottomrule
\end{tabularx}
\end{table}

\section{Recurso de queja por apelaci\'on denegada}

Por medio de este recurso, el \'organo superior resuelve acerca de la decisi\'on del inferior sobre la \textbf{admisibilidad} del recurso. La queja es un \textbf{expediente aut\'onomo}: se inicia y termina en la C\'amara.

Debe interponerse dentro del plazo de \textbf{5 d\'ias} de notificada por ministerio de la ley la resoluci\'on que deneg\'o el recurso. Tiene una doble finalidad: asegurar el derecho de defensa en juicio de las partes y posibilitar el ejercicio de la doble instancia.

\subsection{C\'omo se interpone}

Se debe acompa\~nar \textbf{copia simple, firmada por el letrado del recurrente}, de:
\begin{itemize}
    \item el escrito que dio origen a la resoluci\'on recurrida;
    \item la resoluci\'on recurrida;
    \item el escrito de interposici\'on del recurso;
    \item la resoluci\'on que deneg\'o la apelaci\'on.
\end{itemize}

Adem\'as, en el escrito debe indicarse la fecha de la resoluci\'on recurrida, la fecha de interposici\'on del recurso y la fecha en que se notific\'o la denegatoria.

\subsection{Tr\'amite}

El tr\'amite de la queja \textbf{no suspende} la sentencia ni el tr\'amite del proceso. Decide \'unicamente la C\'amara, sin sustanciaci\'on: no corre traslado a ninguna de las partes. Debe resolver si el recurso fue bien o mal denegado:

\begin{itemize}
    \item si considera que fue \textbf{mal denegado}, ordena al juzgado que disponga su tramitaci\'on;
    \item si considera que fue \textbf{bien denegado}, queda firme la denegatoria.
\end{itemize}

\section{Cuadro comparativo de los recursos}

\begin{table}[H]
\centering
\small
\renewcommand{\arraystretch}{1.3}
\begin{tabularx}{\textwidth}{
    >{\raggedright\arraybackslash}p{0.15\textwidth}
    >{\raggedright\arraybackslash}p{0.14\textwidth}
    >{\raggedright\arraybackslash}p{0.14\textwidth}
    >{\raggedright\arraybackslash}X
    >{\raggedright\arraybackslash}X
}
\toprule
\textbf{Recurso} & \textbf{Se interpone ante} & \textbf{Plazo} & \textbf{Procede contra} & \textbf{Sustanciaci\'on} \\
\midrule
Aclaratoria & El mismo juez & 3 d\'ias & Resoluciones con omisi\'on o error a corregir & No \\
Revocatoria (reposici\'on) & El mismo juez & 3 d\'ias, con apelaci\'on subsidiaria & Providencias simples \'unicamente & No se menciona \\
Apelaci\'on & El mismo juez & 5 d\'ias & Sentencias definitivas de ordinarios, interlocutorias y providencias simples con gravamen irreparable & Seg\'un concesi\'on: libre (agravios) o en relaci\'on (memorial) \\
Queja por apelaci\'on denegada & El \'organo superior (C\'amara) & 5 d\'ias desde la notificaci\'on de la denegatoria & La resoluci\'on que deneg\'o la apelaci\'on & No: la C\'amara decide sin correr traslado \\
Extraordinario (Corte) & Corte Suprema de Justicia de la Naci\'on & 10 d\'ias & Sentencia del tribunal superior, con garant\'ia constitucional afectada & Se da traslado por 10 d\'ias \\
\bottomrule
\end{tabularx}
\end{table}

\section{Recurso extraordinario ante la Corte Suprema}

El recurso extraordinario es dif\'icil y no resulta sencillo. Debe tenerse siempre presente que se interpone cuando est\'an en juego \textbf{garant\'ias constitucionales}: no puede interponerse porque s\'i o porque se lo considere el mejor recurso. Debe considerarse que hay una garant\'ia constitucional afectada.

Se dispone de un plazo de \textbf{10 d\'ias} para interponerlo y fundamentarlo, y siempre debe contarse con la \textbf{sentencia del tribunal superior} que interviene en la causa (en el caso explicado, una sentencia de C\'amara).

\begin{important}
Al recurrir a la Corte debe llenarse un \textbf{formulario}. Si el recurso se interpone sin acompa\~nar el formulario, \textbf{se rechaza}.
\end{important}

Una vez notificado, se da traslado del recurso extraordinario por un plazo de \textbf{10 d\'ias} para contestar, y la notificaci\'on se realiza por c\'edula.

En los recursos extraordinarios \textbf{no pueden copiarse los fundamentos de los recursos anteriores}. Debe especificarse directamente cu\'al es la garant\'ia constitucional afectada y fundamentarla en base a ella. No puede decirse ``como se dijo en la apelaci\'on''. Lo mismo rige al contestar el recurso: no puede decirse ``como se dijo al contestar el memorial de la apelaci\'on'', sino que debe contestarse espec\'ificamente sobre la garant\'ia en juego. De lo contrario, el recurso se rechaza.

\subsection{Consultas sobre plazos y formulario}

Ante la consulta sobre si se trata de dos plazos de 10 d\'ias o de un \'unico plazo de 10 d\'ias dentro del cual deben realizarse la presentaci\'on y el formulario, se explic\'o que, al interponer el recurso, debe acompa\~narse el formulario: todo junto. El formulario debe acompa\~narlo tanto quien inicia (quien interpone) como quien contesta. En consecuencia, casi todo debe realizarse dentro de los 10 d\'ias del extraordinario.

\begin{note}
Estas reglas deben tenerse presentes el d\'ia en que se interponga un recurso extraordinario o el d\'ia en que deba contestarse.
\end{note}

\section{Relaci\'on entre los recursos: ejemplo cotidiano}

El material original ilustra la relaci\'on entre los distintos recursos con el reclamo por la nota de un examen que se considera injusta.

\begin{example}{Reclamar la nota de un examen}
\begin{itemize}
    \item \textbf{Aclaratoria:} se pide que se corrija un error de carga (se coloc\'o mal el nombre o se sum\'o mal).
    \item \textbf{Revocatoria:} se pide al mismo profesor que reconsidere su decisi\'on.
    \item \textbf{Apelaci\'on:} si no cambia su decisi\'on, se acude a una autoridad superior (la direcci\'on de carrera).
    \item \textbf{Queja:} si la direcci\'on ni siquiera permite presentar el reclamo, se reclama contra esa negativa.
    \item \textbf{Recurso extraordinario:} es el \'ultimo paso, solo si se viol\'o una garant\'ia de fondo.
\end{itemize}
Todo ello tiene plazos cortos y reglas de forma. La rebeld\'ia es el caso inverso: se cita a rendir y no se comparece; el tr\'amite sigue sin la persona y cada decisi\'on produce efectos sobre ella aunque no est\'e presente. Ambos bloques son dos caras de la misma regla: en el proceso hay que presentarse y reaccionar a tiempo.
\end{example}

% ------------------------------------------------------------
% CORNELL CAPÍTULO 1
% ------------------------------------------------------------
\section{Cuadro Cornell: los recursos procesales}

\begingroup
\renewcommand{\arraystretch}{1.25}
\begin{longtable}{
    >{\raggedright\arraybackslash}p{0.26\textwidth}
    >{\raggedright\arraybackslash}p{0.66\textwidth}
}
\toprule
\textbf{Preguntas / palabras clave} & \textbf{Notas / respuestas} \\
\midrule
\endfirsthead
\toprule
\textbf{Preguntas / palabras clave} & \textbf{Notas / respuestas} \\
\midrule
\endhead
\midrule
\multicolumn{2}{r}{\small contin\'ua en la p\'agina siguiente} \\
\endfoot
\bottomrule
\endlastfoot

\textbf{¿Qu\'e es un recurso?} &
Medio procesal por el cual la parte agraviada por una resoluci\'on judicial pide su revocaci\'on total o parcial, ante el mismo juez que la dict\'o o ante uno de jerarqu\'ia superior. \\[0.6em]

\textbf{¿C\'omo se clasifican los recursos?} &
Ordinarios (se conceden en cualquier etapa del proceso): aclaratoria, reposici\'on o revocatoria, apelaci\'on y queja por apelaci\'on denegada. Extraordinarios (se conceden excepcionalmente, con procedimiento determinado por la ley): recurso extraordinario ante la Corte Suprema y el del art\'iculo 14 de la Ley 48. \\[0.6em]

\textbf{¿Cu\'ales son los tres puntos b\'asicos para recurrir?} &
Saber qu\'e recurso se interpone, conocer el plazo y saber c\'omo se fundamenta. \\[0.6em]

\textbf{¿C\'omo funciona la aclaratoria?} &
Se interpone ante el mismo juez, en 3 d\'ias, y se fundamenta en el mismo escrito. Sirve para corregir errores u omisiones (una fecha, un nombre). No se sustancia: no se corre traslado a la contraria. \\[0.6em]

\textbf{¿Por qu\'e la revocatoria debe ir con apelaci\'on subsidiaria?} &
La revocatoria procede \'unicamente contra providencias simples, se interpone en 3 d\'ias y se fundamenta en el mismo escrito. Como no suspende el plazo de la apelaci\'on, si se la deniega sin haber apelado en subsidio, el plazo para apelar ya se perdi\'o. \\[0.6em]

\textbf{¿Qu\'e es la apelaci\'on y cu\'ando procede?} &
Se pide al \'organo superior que revoque la decisi\'on del inferior. Se interpone ante el mismo juez en 5 d\'ias. Procede contra sentencias definitivas en juicios ordinarios, interlocutorias y providencias simples que causen gravamen irreparable. \\[0.6em]

\textbf{¿Qu\'e efectos tiene la apelaci\'on?} &
Suspensivo (regla general: se suspende la sentencia) y devolutivo (la sentencia se cumple aunque est\'e apelada: medidas cautelares, alimentos y sumar\'isimos). Tambi\'en existe el efecto diferido, vinculado con el momento de la fundamentaci\'on. \\[0.6em]

\textbf{¿Qu\'e diferencia hay entre concesi\'on libre y en relaci\'on?} &
Libremente: sentencia definitiva de ordinarios; se funda con expresi\'on de agravios ante la C\'amara en 10 d\'ias desde la notificaci\'on de la sala; si no se presenta o no hay cr\'itica razonada y fundada, se tiene por desierto. En relaci\'on: interlocutorias y providencias simples con gravamen irreparable; se funda con memorial ante el mismo juez en 5 d\'ias; traslado a la contraria por 5 d\'ias; luego se pide la elevaci\'on a la C\'amara. \\[0.6em]

\textbf{¿Qu\'e son los incidentes del art\'iculo 250?} &
Se forman cuando la apelaci\'on se concede en relaci\'on, y requieren copias porque la resoluci\'on debe seguir cumpli\'endose. Sentencia definitiva: el expediente va a la C\'amara y las copias quedan en el juzgado. Interlocutoria: el expediente queda y viajan el escrito de interposici\'on, los memoriales y las copias indicadas. \\[0.6em]

\textbf{¿Qu\'e es la queja por apelaci\'on denegada?} &
Expediente aut\'onomo que se inicia y termina en la C\'amara; 5 d\'ias desde la notificaci\'on de la denegatoria por ministerio de la ley. Se acompa\~nan copias simples firmadas por el letrado y se indican las tres fechas. No suspende el proceso ni se sustancia: la C\'amara decide si el recurso fue bien o mal denegado. \\[0.6em]

\textbf{¿Cu\'ando procede el recurso extraordinario y qu\'e exige?} &
Solo cuando hay una garant\'ia constitucional en juego, con sentencia del tribunal superior; 10 d\'ias para interponerlo y fundarlo; formulario obligatorio (sin \'el se rechaza), tanto para quien inicia como para quien contesta; traslado de 10 d\'ias por c\'edula; los fundamentos deben ser propios, sin remitir a recursos anteriores. \\

\midrule
\multicolumn{2}{p{0.94\textwidth}}{
\textbf{Resumen:}
El proceso ofrece herramientas para corregir los errores del juez. Cada recurso responde a tres preguntas: qu\'e se interpone, en qu\'e plazo y con qu\'e fundamentos. La aclaratoria y la revocatoria se resuelven ante el mismo juez en 3 d\'ias; la apelaci\'on, en 5 d\'ias, abre la segunda instancia, con efecto suspensivo o devolutivo y fundamentaci\'on distinta seg\'un se conceda libremente (expresi\'on de agravios, 10 d\'ias) o en relaci\'on (memorial, 5 d\'ias, con incidentes y copias en los casos del art\'iculo 250); la queja controla la denegatoria de la apelaci\'on; y el extraordinario, reservado a garant\'ias constitucionales, exige formulario y fundamentos propios en 10 d\'ias.
} \\
\end{longtable}
\endgroup

% ------------------------------------------------------------
\chapter[Rebeld\'ia]{Rebeld\'ia}
% ------------------------------------------------------------

\section{Concepto y presupuestos}

La rebeld\'ia est\'a regulada a partir del \textbf{art\'iculo 59 del C\'odigo Procesal}.

\begin{definition}{Rebeld\'ia}
Situaci\'on procesal que se configura \textbf{siempre a pedido de parte}: una persona \textbf{debidamente citada} que no comparece en el plazo fijado por el juez (por ejemplo, el de contestaci\'on de la demanda) o que, habiendo comparecido, abandona el proceso.
\end{definition}

Los casos de abandono del proceso se encuentran en los art\'iculos siguientes al 59.

% TODO: contenido incompleto o ambiguo en las notas originales (los art\'iculos sobre abandono del proceso no se identifican con claridad)

\subsection{El \'unico caso en que procede sin pedido de parte}

En los casos de representaci\'on a trav\'es de poder, que se hace mediante un contrato de mandato, puede ocurrir que el abogado renuncie al patrocinio, que el cliente revoque el poder o que alguna de las partes fallezca. Son causales de extinci\'on del contrato de mandato. En ese caso, si no se presenta con nuevo patrocinio letrado, el proceso se cursa en rebeld\'ia.

\begin{important}
Este es el \textbf{\'unico caso} en el cual la rebeld\'ia procede \textbf{sin necesidad de pedido de parte}.
\end{important}

\section{Efectos de la rebeld\'ia: las notificaciones}

Por el art\'iculo 135 ciertas resoluciones deben notificarse por c\'edula (traslado de la demanda, liquidaciones, etc\'etera). Declarada la rebeld\'ia, las notificaciones se hacen \textbf{por ministerio de la ley}, de modo que la c\'edula se utiliza para notificar \'unicamente:

\begin{itemize}
    \item la \textbf{providencia que declara la rebeld\'ia}; y
    \item la \textbf{sentencia}.
\end{itemize}

El resto de las resoluciones, incluidas las liquidaciones, se notifica por ministerio de la ley, en los d\'ias de nota (martes y viernes, o el siguiente d\'ia h\'abil).

Respecto de la prueba confesional, se explic\'o que, estando la rebeld\'ia declarada y firme, se notifica por nota. En el caso de la incomparecencia, por el art\'iculo 41 se notificar\'ia tambi\'en por c\'edula la citaci\'on a absolver posiciones. No obstante, generalmente, si se declar\'o la rebeld\'ia, no se notifica nada m\'as que la declaraci\'on de rebeld\'ia y la sentencia.

% TODO: contenido incompleto o ambiguo en las notas originales (la respuesta sobre la confesional y el art\'iculo 41 es parcialmente ininteligible en el audio y debe contrastarse con el C\'odigo Procesal)

\begin{table}[H]
\centering
\renewcommand{\arraystretch}{1.3}
\begin{tabularx}{\textwidth}{
    >{\raggedright\arraybackslash}p{0.27\textwidth}
    >{\raggedright\arraybackslash}X
    >{\raggedright\arraybackslash}X
}
\toprule
\textbf{Acto} & \textbf{Forma de notificaci\'on} & \textbf{Comentario} \\
\midrule
Providencia que declara la rebeld\'ia & C\'edula & Siempre \\
Sentencia & C\'edula & Siempre \\
Dem\'as resoluciones (incluidas liquidaciones) & Ministerio de la ley (d\'ias de nota: martes y viernes, o el siguiente h\'abil) & Una vez declarada y firme la rebeld\'ia \\
Citaci\'on a absolver posiciones (confesional) & C\'edula (art. 41), seg\'un lo explicado & Como regla general, no se notifica nada m\'as que la declaraci\'on de rebeld\'ia y la sentencia \\
Cuando el rebelde se presenta & C\'edula, en lo que corresponda & Cesa el estado de rebeld\'ia \\
\bottomrule
\end{tabularx}
\end{table}

\section{Efectos de la rebeld\'ia: medidas cautelares y hechos}

\subsection{Medidas cautelares}

La rebeld\'ia, una vez \textbf{declarada y firme} (es decir, notificada en el domicilio real y vencido el plazo para apelar), habilita a trabar medidas cautelares. Adem\'as de los requisitos que deben acreditarse para una medida cautelar (contracautela, peligro en la demora y verosimilitud del derecho), la medida se ve fundada en que la persona, debidamente citada, no quiso comparecer. Ello da un indicio de que, en caso de ser condenada, tampoco querr\'a cumplir; por eso las medidas cautelares proceden.

Si el rebelde es el actor, puede trabarse una medida cautelar para que responda por las \textbf{costas del proceso}. Las medidas cautelares en rebeld\'ia \textbf{subsisten hasta que termine el proceso}.

\subsection{Los hechos: art\'iculo 356 y reconocimiento de la documental}

El art\'iculo 356 impone al demandado el deber de reconocer o negar los hechos. Respecto del rebelde, dice el art\'iculo que los hechos podr\'an ser considerados verdaderos a la luz del juez, seg\'un lo producido en autos. Se trata de una facultad. Con la rebeld\'ia, en caso de duda, esos hechos se reconocen como verdaderos. Lo mismo ocurre con la documental.

En consecuencia, si con la prueba documental pueden acreditarse casi todos los extremos de la pretensi\'on, probablemente se dicte el llamado de autos y no haya mucha prueba que producir. Ello no implica que no se abra a prueba el expediente, porque subsisten las facultades del juez para ahondar en los hechos.

\section{El rebelde que comparece}

Que el rebelde comparezca no implica que el proceso se retrotraiga, porque sigue operando el \textbf{principio de preclusi\'on}. El rebelde puede producir prueba de alg\'un hecho nuevo, asistir a las audiencias, interrogar testigos y, en el caso de una pericia, solicitar que se ampl\'ien los puntos de pericia; puede hacer cualquier cosa, menos aquello que deb\'io hacer al momento de contestar la demanda.

Cuando se presenta, las notificaciones que deben hacerse por c\'edula se hacen por c\'edula, porque cesa el estado de rebeld\'ia.

Respecto de las cautelares, podr\'a pedir su \textbf{levantamiento} siempre que acredite que la rebeld\'ia fue una circunstancia ajena a su voluntad y que no pudo evitarla (por ejemplo, no estaba en el pa\'is o estaba en coma, de modo que no pudo presentarse en el proceso). Se aplican, en este caso, las medidas para ampliar, reducir o sustituir las medidas cautelares.

\begin{important}
El proceso no se retrotrae si el rebelde comparece (preclusi\'on): toma el proceso en el estado en que est\'a. Para levantar las cautelares debe probar que la rebeld\'ia fue una circunstancia ajena a su voluntad que no pudo evitar.
\end{important}

\section{Nulidad, prueba en segunda instancia y alimentos}

\subsection{Prueba en segunda instancia}

Si el rebelde se presenta antes del dictado de la sentencia y apela, en los t\'erminos del \textbf{art\'iculo 260, inciso 5, apartado a)}, puede solicitar que se abra a prueba en segunda instancia para producir la prueba que se vio impedido de producir en primera. Esto deriva del principio de preclusi\'on procesal. Queda sujeto a que la notificaci\'on cursada haya sido v\'alida (dirigida al domicilio real y dem\'as requisitos).

Si se presenta antes del dictado de la sentencia, cesa el estado de rebeld\'ia y las notificaciones vuelven a ser las de un expediente normal, de modo que puede apelar la sentencia. La prueba se produce en segunda instancia, por el art\'iculo 260: es la prueba que se vio impedido de producir en primera. Para la imposici\'on de las costas se tienen en cuenta los perjuicios que la rebeld\'ia caus\'o en el proceso.

\subsection{Nulidad de la notificaci\'on}

Si se anula la notificaci\'on a trav\'es de un incidente, cae la declaraci\'on de rebeld\'ia y, con ella, \textbf{todo el proceso}.

\begin{example}{Notificaci\'on a un domicilio que no era el real}
Si una parte env\'ia una c\'edula a un domicilio y le informan ``no vive ac\'a'', y no adopta las diligencias necesarias ni realiza investigaci\'on alguna y a\'un as\'i env\'ia all\'i la c\'edula, la rebeld\'ia no refleja la realidad: la persona nunca estuvo realmente notificada. Hay que probarlo, pero si se prueba en un incidente de nulidad, se cae la rebeld\'ia y se cae todo el proceso.
\end{example}

Por el \textbf{art\'iculo 67}, ejecutoriada la sentencia, no se admite recurso alguno contra la sentencia dictada en rebeld\'ia. Sin embargo, nada impide plantear un incidente de nulidad directamente de la notificaci\'on y arrastrar todo el proceso en nulidad.

\subsection{La rebeld\'ia en los procesos de alimentos}

En los procesos de alimentos no se declara la rebeld\'ia, porque la ley ya impone la sanci\'on: se dicta sentencia con las constancias que hay en la causa. Adem\'as, cuando se decreta la rebeld\'ia, en todo proceso se tienen por ciertos los dichos de la parte actora, porque no contest\'o.

% TODO: contenido incompleto o ambiguo en las notas originales (el fragmento sobre alimentos es parcialmente ininteligible en el audio; la referencia al derecho del ni\~no se transcribe con la lectura m\'as probable)

\begin{note}
Seg\'un las notas de edici\'on del material original, la respuesta sobre la confesional y el art\'iculo 41, y el fragmento sobre alimentos, son parcialmente ininteligibles en el audio y se transcribieron con la lectura m\'as probable. Se recomienda contrastarlos con el texto del C\'odigo Procesal.
\end{note}

% ------------------------------------------------------------
% CORNELL CAPÍTULO 2
% ------------------------------------------------------------
\section{Cuadro Cornell: rebeld\'ia}

\begingroup
\renewcommand{\arraystretch}{1.25}
\begin{longtable}{
    >{\raggedright\arraybackslash}p{0.26\textwidth}
    >{\raggedright\arraybackslash}p{0.66\textwidth}
}
\toprule
\textbf{Preguntas / palabras clave} & \textbf{Notas / respuestas} \\
\midrule
\endfirsthead
\toprule
\textbf{Preguntas / palabras clave} & \textbf{Notas / respuestas} \\
\midrule
\endhead
\midrule
\multicolumn{2}{r}{\small contin\'ua en la p\'agina siguiente} \\
\endfoot
\bottomrule
\endlastfoot

\textbf{¿Qu\'e es la rebeld\'ia y cu\'ando procede?} &
Art. 59 del C\'odigo Procesal: persona debidamente citada que no comparece en el plazo fijado por el juez o que abandona el proceso habiendo comparecido. Se configura siempre a pedido de parte. \\[0.6em]

\textbf{¿Cu\'al es el \'unico caso en que procede sin pedido de parte?} &
Cuando el mandato se extingue (renuncia al patrocinio, revocaci\'on del poder, fallecimiento de una de las partes) y no se presenta nuevo patrocinio letrado. \\[0.6em]

\textbf{¿C\'omo se notifica a quien es rebelde?} &
Por ministerio de la ley. Se usa c\'edula solo para la providencia que declara la rebeld\'ia y para la sentencia; el resto, por nota (martes y viernes, o el siguiente d\'ia h\'abil). Sobre la confesional se cita el art\'iculo 41 (punto a contrastar con el C\'odigo). \\[0.6em]

\textbf{¿Qu\'e efectos tiene sobre las medidas cautelares?} &
Declarada y firme, habilita las cautelares, adem\'as de contracautela, peligro en la demora y verosimilitud del derecho, porque la falta de comparecencia es un indicio de que no se cumplir\'a la condena. Si el rebelde es el actor, puede trabarse una cautelar para que responda por las costas. Subsisten hasta que termine el proceso. \\[0.6em]

\textbf{¿Qu\'e efectos tiene sobre los hechos y la documental?} &
Por el art\'iculo 356, los hechos pueden ser considerados verdaderos seg\'un lo producido en autos (facultad del juez); con la rebeld\'ia, en caso de duda, se reconocen como verdaderos, y lo mismo ocurre con la documental. El juez conserva facultades para ahondar en los hechos. \\[0.6em]

\textbf{¿Qu\'e puede hacer el rebelde que comparece?} &
El proceso no se retrotrae (preclusi\'on). Puede producir prueba de hechos nuevos, ir a audiencias, interrogar testigos y pedir ampliaci\'on de puntos de pericia, pero no lo que deb\'io hacer al contestar la demanda. Puede pedir el levantamiento de cautelares si prueba que la rebeld\'ia le fue ajena y no pudo evitarla. \\[0.6em]

\textbf{¿Cu\'ando puede producir prueba en segunda instancia?} &
Si se presenta antes de la sentencia y apela (art. 260, inc. 5, ap. a), puede pedir apertura a prueba en segunda instancia para producir la prueba que no pudo producir en primera, siempre que la notificaci\'on haya sido v\'alida. \\[0.6em]

\textbf{¿Cu\'ando cae todo el proceso?} &
Si se anula la notificaci\'on mediante incidente de nulidad, cae la declaraci\'on de rebeld\'ia y, con ella, todo el proceso. El art\'iculo 67 no admite recurso contra la sentencia dictada en rebeld\'ia una vez ejecutoriada, pero nada impide plantear la nulidad de la notificaci\'on. \\[0.6em]

\textbf{¿Qu\'e ocurre con la rebeld\'ia en alimentos?} &
No se declara rebeld\'ia: la ley ya impone la sanci\'on y se dicta sentencia con las constancias de la causa. En todo proceso con rebeld\'ia se tienen por ciertos los dichos de la parte actora. \\

\midrule
\multicolumn{2}{p{0.94\textwidth}}{
\textbf{Resumen:}
La rebeld\'ia, regulada desde el art\'iculo 59, sanciona la inacci\'on: altera el r\'egimen de notificaciones, habilita medidas cautelares y permite tener por ciertos los hechos. Quien comparece puede continuar el proceso sin retrotraerlo y, si la notificaci\'on fue inv\'alida, plantear la nulidad que derriba todo lo actuado. Ambos temas confluyen en una ense\~nanza pr\'actica: los plazos son breves y presentarse y fundar a tiempo es lo que protege el derecho de defensa.
} \\
\end{longtable}
\endgroup

\end{document}
